% Part: incompleteness
% Chapter: introduction
% Section: overview

\documentclass[../../../include/open-logic-section]{subfiles}

\begin{document}

\olfileid[hi]{inc}{int}{ovr}

\olsection{अपूर्णता परिणामों का अवलोकन}

हिल्बर्ट को आशा थी कि गणित को किसी ऐसे !!{axiomatizable} सिद्धांत में
औपचारिक बनाया जा सकता है जिसे !!{complete} और !!{decidable} सिद्ध करना संभव
होगा। इसके अतिरिक्त, उनका लक्ष्य बहुत दुर्बल, ``परिमितवादी'' साधनों से इस
सिद्धांत की संगतता सिद्ध करना था, जिससे अंतःप्रज्ञावाद की चुनौतियों के विरुद्ध
चिरसम्मत गणित की रक्षा हो सके। ग्योडेल (G\"odel) के अपूर्णता प्रमेयों ने
दिखाया कि ये लक्ष्य प्राप्त नहीं किए जा सकते।

ग्योडेल (G\"odel) के प्रथम अपूर्णता प्रमेय ने दिखाया कि रसेल और व्हाइटहेड की
\emph{प्रिंसिपिया मैथमेटिका} का एक रूप !!{complete} नहीं है। किंतु प्रमाण
वास्तव में अत्यंत सामान्य था और अनेक प्रकार के सिद्धांतों पर लागू होता है।
इसका अर्थ है कि केवल \emph{प्रिंसिपिया मैथमेटिका} ही गणित को पूर्णतः ग्रहण
करने में विफल नहीं हुई, बल्कि \emph{कोई भी} स्वीकार्य सिद्धांत ऐसा नहीं कर
सकता। उन सिद्धांतगत विशेषताओं को अलग करने में कुछ समय लगा जो अपूर्णता
प्रमेयों के लागू होने के लिए पर्याप्त हैं, और ग्योडेल (G\"odel) के प्रमाण को
इस प्रकार सामान्य बनाने में भी कि वह केवल इन्हीं विशेषताओं पर निर्भर रहे।
किंतु अब हम अंकगणित की भाषा $\Lang L_A$ के सिद्धांतों के लिए प्रथम अपूर्णता
प्रमेय का अत्यंत सामान्य रूप कह सकते हैं।

\begin{thm}
यदि $\Gamma$,~$\Lang L_A$ में ऐसा संगत और !!{axiomatizable} सिद्धांत है जो
सभी संगणनीय फलनों और !!{decidable} संबंधों को !!{represents} करता है, तो
$\Gamma$ !!{complete} नहीं है।
\end{thm}

$\Gamma$ के !!{complete} न होने का अर्थ है कि कम-से-कम एक !!{sentence}~$!A$
के लिए $\Gamma \Proves/ !A$ और $\Gamma \Proves/ \lnot !A$। ऐसे
!!a{sentence} को~($\Gamma)$ से \emph{स्वतंत्र} कहते हैं। वास्तव में हम अपेक्षाकृत
शीघ्र सिद्ध कर सकते हैं कि स्वतंत्र वाक्य अवश्य होते हैं। किंतु प्रमेय के
ग्योडेल (G\"odel) वाले प्रमाण की शक्ति इस तथ्य में है कि वह ऐसे स्वतंत्र
!!{sentence} का एक \emph{विशिष्ट उदाहरण} प्रस्तुत करता है। यह रोचक रचना
!!a{sentence}~$!G_\Gamma$ उत्पन्न करती है, जिसे~$\Gamma$ का \emph{ग्योडेल (G\"odel)
 वाक्य} कहते हैं। वह असिद्धनीय है क्योंकि $\Gamma$ में $!G_\Gamma$
इस दावे के तुल्य है कि~$\Gamma$ में $!G_\Gamma$ असिद्धनीय है। रचना यह काम
\emph{रचनात्मक रूप से} करती है; अर्थात्~$\Gamma$ का अभिगृहीतीकरण और
!!{derivation} तंत्र का वर्णन दिए जाने पर प्रमाण वास्तव में~$!G_\Gamma$ लिखने
की एक विधि देता है।

ग्योडेल (G\"odel) के प्रमाण की रचना के लिए हमें~$\Lang L_A$ में स्वयं
$\Lang L_A$ के पदों और !!{formula} के गुण तथा उन पर संक्रियाएँ व्यक्त करने
का कोई उपाय खोजना होगा। इनमें ``$!A$ एक !!a{sentence} है'', ``$\delta$,
~$!A$ की !!a{derivation} है'' जैसे गुण और $\Subst{!A}{t}{x}$ जैसी संक्रियाएँ
सम्मिलित हैं। इस उपाय को (क) चिह्नों और उनके अनुक्रमों (जिनसे पद, !!{formula},
!!{derivation} आदि बनते हैं) के प्राकृतिक संख्याओं के रूप में ``कूटलेखन'' के
माध्यम से इन गुणों और संबंधों को व्यक्त करना होगा (क्योंकि $\Lang L_A$
प्राकृतिक संख्याओं की ही चर्चा कर सकती है)। उसे (ख) यह ऐसे ढंग से करना होगा
कि $\Gamma$ संबंधित तथ्यों को सिद्ध करे; इसलिए हमें दिखाना होगा कि इन गुणों का
कूटलेखन प्राकृतिक संख्याओं के !!{decidable} गुणों से होता है और संक्रियाएँ
प्राकृतिक संख्याओं पर संगणनीय फलनों के अनुरूप हैं। इसे ``वाक्यविन्यास का
अंकगणितीकरण'' कहते हैं।

किंतु वाक्यविन्यास का अंकगणितीकरण कैसे हो सकता है, इसकी जाँच से पहले हम इस
शर्त पर विचार करेंगे कि~$\Gamma$ ``पर्याप्त प्रबल'' है, अर्थात् सभी संगणनीय
फलनों और !!{decidable} संबंधों को !!{represents} करता है। इसके लिए हमें
``संगणनीय'' की सटीक परिभाषा देनी होगी। यह कई प्रकार से किया जा सकता है;
उदाहरणतः ट्यूरिंग मशीनों के प्रतिरूप द्वारा, अथवा किसी सामान्य-प्रयोजन
क्रमादेश भाषा के क्रमादेशों से संगणनीय फलनों के रूप में। किंतु हमारा लक्ष्य
इन फलनों और संबंधों को भाषा~$\Lang L_A$ के किसी सिद्धांत में !!{represents}
करना है, इसलिए केवल संख्यात्मक फलनों की संगणनीयता की कोई सरल परिभाषा चुनना
उत्तम है। यह \emph{पुनरावर्ती फलन} की धारणा है। अतः हम पहले पुनरावर्ती फलनों
की चर्चा करेंगे। फिर दिखाएँगे कि $\Th{Q}$ पहले से ही सभी पुनरावर्ती फलनों और
संबंधों को !!{represents} करता है। इससे हम अपूर्णता प्रमेय को $\Th{Q}$ और
$\Th{PA}$ जैसे विशिष्ट सिद्धांतों पर लागू कर सकेंगे, क्योंकि हम स्थापित कर
चुके होंगे कि ये ``पर्याप्त प्रबल'' सिद्धांतों के उदाहरण हैं।

वाक्यविन्यास के अंकगणितीकरण का अंतिम परिणाम !!a{formula}
$\OProv[\Gamma](x)$ है, जो !!{formula} को संख्याओं के रूप में कूटित करके
~$\Gamma$ के अभिगृहीतों से सिद्धनीयता व्यक्त करता है। विशेषतः, यदि $!A$ को
संख्या~$n$ से कूटित किया गया है और $\Gamma \Proves !A$, तो $\Gamma \Proves
\Prov[\Gamma](\num{n})$। $\Gamma$ के लिए यह ``सिद्धनीयता विधेय'' हमें एक
निश्चित अर्थ में $\Gamma$ की संगतता को भी~$\Lang L_A$ के !!a{sentence} के रूप
में व्यक्त करने देता है:~$\Gamma$ का ``संगतता कथन'' वह !!{sentence} $\lnot
\Prov[\Gamma](\num{n})$ हो, जहाँ हम $n$ को किसी विरोधाभास, उदाहरणतः~$\lfalse$,
का कूट लेते हैं। द्वितीय अपूर्णता प्रमेय कहता है कि संगत !!{axiomatizable}
सिद्धांत अपने ही संगतता कथनों को भी सिद्ध नहीं करते। इस प्रमेय के लागू होने
के लिए आवश्यक शर्तें केवल यह कहने से कुछ अधिक कठोर हैं कि सिद्धांत सभी
संगणनीय फलनों और !!{decidable} संबंधों को निरूपित करता है; किंतु हम दिखाएँगे
कि $\Th{PA}$ उन्हें पूरा करता है।

\end{document}
